\documentclass[10pt,a4paper,twocolumn]{article}

\usepackage[utf8]{inputenc}
\usepackage[T1]{fontenc}
\usepackage[brazilian]{babel}
\usepackage{lmodern}
\usepackage{microtype}
\usepackage[margin=1.8cm,columnsep=0.7cm]{geometry}
\usepackage{amsmath,amssymb}
\usepackage{booktabs}
\usepackage{graphicx}
\usepackage{enumitem}
\usepackage[font=small,labelfont=bf]{caption}
\usepackage[hidelinks]{hyperref}

\setlist{nosep,leftmargin=1.5em}
\setlength{\parindent}{1em}
\renewcommand{\arraystretch}{1.1}

% Figura com marcador de lugar caso o arquivo nao exista em figuras/
\newcommand{\fig}[2]{%
  \IfFileExists{#1}{\includegraphics[width=#2\linewidth]{#1}}%
  {\fbox{\parbox{#2\linewidth}{\centering\small\vspace{1.2cm}[figura: \texttt{\detokenize{#1}}]\vspace{1.2cm}}}}}

\newcommand{\F}{\texttt{FACCAO}}

\title{\textbf{Do Simbólico à Incerteza: inferindo a presença e a expansão
de facções criminosas nos municípios da Amazônia Legal}}
\author{Equipe 1 --- P3 Segurança Pública\\[2pt]
\small Gabriel Menezes Rodrigues, Luiz Felipe Gonzaga do Carmo, Joshua Tadeu Bastos Pereira,\\
\small Romário Roberto de Pádua Andrade, Ruan Pablo Nascimento de Aragão, Fernando Teles dos Reis\\[2pt]
\small Fundamentos de Inteligência Artificial --- Prof.\ Edjard Mota --- IComp/UFAM, 2026/2}
\date{Novembro de 2026}

\begin{document}

\maketitle
\begin{abstract}
\noindent A presença de facções criminosas não aparece em nenhuma base oficial: é uma
variável oculta. Combinamos as estatísticas de violência do SINESP VDE (2020--2024) e a
população do IBGE para inferi-la nos 808 municípios dos estados da Amazônia Legal, usando
como referência o \emph{Cartografias da Violência na Amazônia} de 2024. Percorremos as
seis etapas do pipeline: árvore de decisão, modelos lineares, probabilidade, Naive Bayes
com EM, rede bayesiana e lógica fuzzy. A rede bayesiana diagnóstica atinge AUC de 0,81
(0,80 testando num estado que não viu). Sem usar nenhum dado de 2025, a probabilidade
calculada antecipa quais municípios passaram a ter facção segundo a edição de 2025
(AUC 0,76): na prioridade fuzzy baixa, 13\% entraram; na alta, 44\%. A violência indica
onde a facção tende a estar, mas não quando ela chega.
\end{abstract}

%=====================================================================
\section{Introdução}

Nas duas últimas décadas, facções criminosas nascidas no Sudeste (PCC e Comando Vermelho)
e grupos regionais se espalharam pela Amazônia, atraídos pelas rotas fluviais e terrestres
da cocaína produzida na Bolívia, no Peru e na Colômbia \cite{r11} e pelo sistema prisional,
onde se organizam \cite{r10}. A literatura descreve uma interiorização das facções na
região, ligada também ao garimpo ilegal \cite{r12,r13}. O \emph{Cartografias da Violência
na Amazônia}, do Fórum Brasileiro de Segurança Pública com o Instituto Mãe Crioula, mapeia
município a município quais facções atuam na região: na 3ª edição (dados de 2024) havia
registro em 259 dos 808 municípios dos nove estados; na 4ª edição (2025), em 344.

Esse mapeamento depende de trabalho de campo e de fontes policiais e não é atualizado
continuamente. As estatísticas de violência, por outro lado, são publicadas todo ano pelo
Ministério da Justiça. Trabalhos próximos usam aprendizado de máquina para prever
homicídios em Belém a partir de boletins de ocorrência \cite{r16} e para prever a
infiltração mafiosa em municípios italianos com validação fora do tempo \cite{r17}; neste
último, os ``falsos positivos'' se parecem com municípios infiltrados ainda não detectados,
ponto que retomamos na Seção~\ref{sec:concl}. A pergunta deste trabalho é:

\begin{quote}\itshape
O município da Amazônia Legal tem presença de facção criminosa (2024), dado o seu porte e
o padrão de violência registrado entre 2021 e 2023? E o modelo antecipa quais municípios
passaram a ter facção em 2025?
\end{quote}

A resposta final é uma prioridade de monitoramento por município, produzida por um sistema
fuzzy que recebe a probabilidade calculada pela rede bayesiana. A prioridade não é uma
acusação: indica onde a inteligência policial teria mais chance de encontrar atuação de
facção.

%=====================================================================
\section{Dados e protocolo de validação}

\paragraph{Fontes.} (i) Rótulos: Cartografias, 3ª edição \cite{r4} (Tabelas 13--21,
presença em 2024) e 4ª edição \cite{r5} (quadros por estado, extraídos do PDF com
\texttt{extrair\_cartografias\_2025.py}). (ii) Atributos: Banco VDE do SINESP/MJSP
\cite{r6}, 2020--2024, por município (homicídio doloso, tentativa, morte a esclarecer,
latrocínio, mandados). (iii) População: IBGE, Censo 2022 e estimativas \cite{r7}, unida
pelo código do município. Os 5.570 municípios do VDE casaram com o IBGE; os 808 dos 9
estados casaram com o Cartografias. (iv) Bases externas, usadas na análise do garimpo e nos
experimentos (Seção~\ref{sec:exp}): autos de infração e termos de destruição do IBAMA
\cite{r8}, alertas do DETER/INPE \cite{r9}, faixa de fronteira, malha municipal e REGIC do
IBGE \cite{r7}, todos de 2021--23 e agregados por município
(\texttt{montar\_bases\_externas.py}). Dos arquivos do IBAMA, que identificam os autuados,
guardamos só contagens por município.

\paragraph{Atributos.} Para a árvore, o Naive Bayes e a rede, usamos atributos discretos e
interpretáveis (Tabela~\ref{tab:attr}). As taxas são médias anuais de 2021--23 por 100 mil
habitantes (Censo 2022).

\begin{table}[t]
\centering\small
\begin{tabular}{@{}lp{5.6cm}@{}}
\toprule
Variável & Definição\\
\midrule
\texttt{PORTE} & pequeno ($<$10 mil), médio (10--50 mil), grande (50 mil+)\\
\texttt{TXH} & taxa de homicídio 2021--23, em tercis\\
\texttt{TXT} & taxa de tentativas de homicídio 2021--23, em tercis\\
\texttt{ESC} & teve morte a esclarecer em 2021--23\\
\texttt{LATRO} & teve latrocínio em 2021--23\\
\texttt{UF/FRONT} & estado / estado faz fronteira com outro país\\
\texttt{ALVO} & facção registrada em 2024 (32,1\%)\\
\bottomrule
\end{tabular}
\caption{Atributos e alvo.}\label{tab:attr}
\end{table}

\paragraph{Validação.} Sempre com \texttt{seed = 42}: (a) validação cruzada estratificada
de 5 dobras, recalculando os tercis de \texttt{TXH} e \texttt{TXT} com os cortes do treino
de cada dobra, para o teste não influenciar a discretização; (b) \emph{deixando um estado
de fora} (treina com 8, testa no 9º); (c) \emph{fora do tempo}: entre os 549 municípios sem
facção em 2024, quais passaram a ter em 2025, rótulo que nenhum modelo viu. Como as classes
são desbalanceadas (``sempre não'' acerta 67,9\%), reportamos também acurácia balanceada e
AUC.

%=====================================================================
\section{Etapa 1 --- Árvore de decisão}

\paragraph{Entropia e ganho (cálculo manual).} Na raiz, 259 dos 808 municípios têm facção:
$H=-\frac{259}{808}\log_2\frac{259}{808}-\frac{549}{808}\log_2\frac{549}{808}=0{,}905$
bits. Para \texttt{PORTE}, os filhos são pequeno (27/267 com facção, $H=0{,}473$), médio
(145/435, $H=0{,}918$) e grande (87/106, $H=0{,}678$). A entropia ponderada é
$(267\cdot0{,}473+435\cdot0{,}918+106\cdot0{,}678)/808=0{,}740$, e o ganho é
$0{,}905-0{,}740=0{,}165$, o maior entre os seis atributos (Tabela~\ref{tab:ganho}).

\begin{table}[t]
\centering\small
\begin{tabular}{@{}lcccccc@{}}
\toprule
& \texttt{PORTE} & \texttt{UF} & \texttt{LATRO} & \texttt{ESC} & \texttt{TXH} & \texttt{TXT}\\
\midrule
Ganho & \textbf{0,165} & 0,151 & 0,070 & 0,030 & 0,027 & 0,017\\
\bottomrule
\end{tabular}
\caption{Ganho de informação na raiz (nível 1).}\label{tab:ganho}
\end{table}

No \textbf{nível 2}, em todos os três ramos de porte o melhor atributo é \texttt{UF}: ganho
de 0,167 entre os pequenos ($H=0{,}473$), 0,133 entre os médios ($H=0{,}918$) e 0,044 entre
os grandes ($H=0{,}678$). No \textbf{nível 3}, os atributos de violência finalmente
aparecem: nos médios do Mato Grosso ($n=61$, $H=0{,}929$) o melhor é \texttt{TXT} (ganho
0,135); nos médios de Rondônia ($n=23$, $H=0{,}999$) e do Amazonas ($n=48$, $H=0{,}696$),
\texttt{TXH} (0,164 e 0,107); nos médios do Maranhão ($n=153$) e do Pará ($n=90$),
\texttt{LATRO} (0,040 e 0,035). A árvore diz, portanto, que o tamanho do município e o
estado explicam a maior parte da presença de facção; a violência refina a decisão dentro de
cada estado.

\paragraph{Implementação e validação.} Implementamos o \textsc{Learn-Decision-Tree} (ID3)
do zero e comparamos com o \texttt{DecisionTreeClassifier} (critério entropia) em 5 dobras
(Tabela~\ref{tab:arvore}). As duas implementações sem poda dão o mesmo resultado (0,749 e
0,745). A poda por profundidade reduz o sobreajuste: o treino cai de 87\% para 82\% e o
teste sobe para 78,6\%. A Figura~\ref{fig:prof} mostra a curva de profundidade.

\begin{table}[t]
\centering\small
\begin{tabular}{@{}lccc@{}}
\toprule
Modelo & Treino & Teste & Bal.\\
\midrule
ID3 própria, sem poda & 0,874 & 0,749 & 0,687\\
ID3 própria, prof.\ 3 & 0,824 & \textbf{0,786} & \textbf{0,708}\\
sklearn, sem poda & 0,875 & 0,745 & 0,683\\
sklearn, prof.\ 4 & 0,799 & 0,782 & 0,679\\
\bottomrule
\end{tabular}
\caption{Validação cruzada 5-fold (tercis calculados no treino de cada
dobra).}\label{tab:arvore}
\end{table}

\begin{figure}[t]
\centering
\fig{figuras/fig2_profundidade.pdf}{0.95}
\caption{Acurácia de teste por profundidade máxima.}\label{fig:prof}
\end{figure}

\paragraph{Monotonicidade.} A regra mais simples sugerida pelos dados é
$\forall x\,(\textit{Pequeno}(x)\rightarrow\neg\textit{Facção}(x))$. Ela cobre 267
municípios e acerta 89,9\%. Mas, ao acrescentarmos um único fato novo, ``o município fica
no Acre'' ou ``em Roraima'', a conclusão se inverte: todos os pequenos desses estados têm
facção (3 de 3 e 2 de 2), e em Rondônia 9 de 22. São estados de fronteira, na rota da
cocaína, onde a facção chega até às cidades pequenas. Na lógica clássica, o raciocínio é
monotônico: se a regra foi aceita como axioma, nenhum fato novo a derruba, e a base passaria
a ser inconsistente. No mundo real, a informação nova invalida a conclusão anterior. Esse é
o limite das regras rígidas que motiva as etapas seguintes, em que as conclusões são graus
de crença revisáveis.

%=====================================================================
\section{Etapa 2 --- Modelos lineares}

\paragraph{Regressão.} Como alvo contínuo usamos $y=\log(1+\text{homicídios em 2023})$. Os
16 atributos são $\log(\text{população})$ e as contagens de 2020--2022 (em $\log(1+x)$).
Divisão 80/20. A regressão univariada por fórmula fechada sobre os homicídios de 2022 dá
$\hat y=0{,}282+0{,}812\,x$ ($R^2=0{,}672$). A multivariada com gradiente descendente do
zero (atributos padronizados, $\alpha=0{,}05$, 5.000 iterações) chega a $R^2=0{,}763$. Como
os 16 atributos são muito correlacionados (o mesmo evento em três anos seguidos),
aplicamos Ridge: o $\lambda$ escolhido por validação cruzada é 0,084, com $R^2=0{,}762$
(Figura~\ref{fig:ridge}). À medida que $\lambda$ cresce, os coeficientes encolhem e deixam
de oscilar entre anos correlacionados.

\begin{figure*}[t]
\centering
\fig{figuras/fig1_pipeline.pdf}{0.95}
\caption{Pipeline integrado. As setas tracejadas indicam reutilização: os atributos e a
discretização definidos na Etapa 1 alimentam a rede, e o EM mostra que o volume de violência
importa, o que orienta a escolha das entradas do sistema fuzzy.}\label{fig:pipeline}
\end{figure*}

\begin{figure}[t]
\centering
\fig{figuras/fig3_ridge.pdf}{0.95}
\caption{Coeficientes do Ridge em função de $\lambda$.}\label{fig:ridge}
\end{figure}

\paragraph{Logística.} Para o problema da Etapa 1, implementamos a regressão logística
(sigmoide + gradiente, com pesos de classe balanceados). Acurácia de 75,9\% e balanceada de
0,737, idênticas ao \texttt{LogisticRegression}. Todos os modelos próprios coincidem com o
scikit-learn até a terceira casa (Tabela~\ref{tab:lin}).

\begin{table}[t]
\centering\small
\begin{tabular}{@{}lcc@{}}
\toprule
Modelo & $R^2$ & Acurácia (bal.)\\
\midrule
Simples (hom.\ 2022 $\to$ 2023) & 0,672 & --\\
Multivariada (GD, 16 atributos) & 0,763 & --\\
Ridge ($\lambda=0{,}084$) & 0,762 & --\\
Logística (facção) & -- & 0,759 (0,737)\\
\bottomrule
\end{tabular}
\caption{Resultados no teste (20\%). O scikit-learn dá os mesmos valores.}\label{tab:lin}
\end{table}

%=====================================================================
\section{Etapa 3 --- Probabilidade}

Cada município é um ``mundo''. Definimos $F$ = tem facção em 2024; $V$ = taxa de homicídio
no tercil mais alto; $E$ = teve morte a esclarecer. A conjunta de $F$ e $V$ está na
Tabela~\ref{tab:conj}.

\begin{table}[t]
\centering\small
\begin{tabular}{@{}lccc@{}}
\toprule
& $F=0$ & $F=1$ & $P(V)$\\
\midrule
$V=0$ & 0,4876 & 0,1795 & 0,6671\\
$V=1$ & 0,1918 & 0,1411 & 0,3329\\
\midrule
$P(F)$ & 0,6795 & 0,3205 & 1,0000\\
\bottomrule
\end{tabular}
\caption{Distribuição conjunta (4 eventos atômicos).}\label{tab:conj}
\end{table}

\paragraph{Bayes: causal $\times$ diagnóstico.} O sentido \emph{causal} é $P(V\mid F)$: a
facção provoca disputas e execuções. Nos dados, $P(V{=}1\mid F{=}1)=0{,}440$. O analista,
porém, observa o efeito e quer a causa oculta, o sentido \emph{diagnóstico}:
\[
P(F\mid V)=\frac{P(V\mid F)\,P(F)}{P(V)}=\frac{0{,}440\times0{,}3205}{0{,}3329}=0{,}424,
\]
contra $P(F\mid\neg V)=0{,}269$: a taxa alta multiplica a chance por 1,6. A conferência
direta na tabela dá o mesmo valor.

\paragraph{Regra da cadeia.} $P(F,V,E)=P(F)\,P(V\mid F)\,P(E\mid F,V)$ reproduz exatamente
as 8 probabilidades conjuntas (soma 1). \textbf{Independência condicional.} Testamos se $V$
e $E$ são independentes dado $F$: para $F=0$, $P(V,E\mid F)=0{,}204$ contra o produto 0,167
($\chi^2=14{,}5$, $p<10^{-3}$); para $F=1$, 0,386 contra 0,350 ($p=0{,}006$). A hipótese do
Naive Bayes é violada: os dois sintomas compartilham uma causa que não está no modelo (o
volume de violência).

\paragraph{Por que a LPO não diz ``42\% de chance''.} A fórmula
$\forall x\,(\textit{ViolênciaAlta}(x)\rightarrow\textit{Facção}(x))$ é falsa, porque há
municípios violentos sem facção; a sua negação também não serve. A lógica de primeira ordem
só atribui verdadeiro ou falso; para dizer ``em 42\% dos casos'' é preciso uma medida de
crença, que é a probabilidade.

\paragraph{Transições 2024 $\to$ 2025.} Com a 4ª edição, a distribuição conjunta dos regimes
nos dois anos (Tabela~\ref{tab:trans}) mostra a dinâmica. Dos 549 municípios sem registro em
2024, 129 (23,5\%) passaram a ter facção em 2025, e a chance cresce com o porte: 13,8\% nos
pequenos, 30,7\% nos médios e 36,8\% nos grandes. Por Bayes,
$P(\text{entrada}\mid\text{grande})=0{,}054\times0{,}235/0{,}035=0{,}368$. Também há 45
``saídas'', o que sugere que parte da mudança vem do método, como o próprio relatório
admite.

\begin{table}[t]
\centering\small
\begin{tabular}{@{}lcccc@{}}
\toprule
2024 $\backslash$ 2025 & sem reg. & uma & disputa & total\\
\midrule
sem registro & 420 & 111 & 18 & 549\\
uma facção & 32 & 126 & 16 & 174\\
disputa & 13 & 18 & 54 & 85\\
\bottomrule
\end{tabular}
\caption{Regimes de facção em 2024 e 2025 (municípios).}\label{tab:trans}
\end{table}

\paragraph{Garimpo e a entrada em 2025.} O garimpo ilegal é descrito como uma das economias
exploradas pelas facções na região, no que se convencionou chamar de ``narcogarimpo''
\cite{r13,r14}. Medimos o garimpo em 2021--23 por duas fontes independentes: auto de
infração do IBAMA por garimpo ou mineração ilegal (71 municípios) e alerta do DETER/INPE da
classe mineração (76). Entre os 549 sem facção em 2024 (Tabela~\ref{tab:garimpo}), os que
tinham garimpo tiveram o dobro de entradas em 2025 nas duas fontes (teste exato de Fisher,
$p=0{,}001$ e $0{,}003$). Por Bayes, com a fonte IBAMA,
$P(\text{entrada}\mid G)=P(G\mid\text{entrada})P(\text{entrada})/P(G)=0{,}147\times0{,}235/0{,}077=0{,}45$.
Já a associação com a facção \emph{existente} em 2024 é fraca: pelo IBAMA, dentro do mesmo
porte, não há diferença (médios: 37\% contra 33\%). O garimpo sinaliza \emph{para onde} a
facção vai, o que a violência não capta (Seção~\ref{sec:fuzzy}). Cuidados: autuação também
reflete onde houve fiscalização, e o DETER inclui mineração legal.

\begin{table}[t]
\centering\small
\begin{tabular}{@{}lcc@{}}
\toprule
Entrada em 2025 & IBAMA & INPE\\
\midrule
Com garimpo & 45,2\% & 44,7\%\\
Sem garimpo & 21,7\% & 21,9\%\\
Médios com garimpo & 55,6\% ($n=27$) & 60,0\% ($n=25$)\\
Médios sem garimpo & 28,1\% & 27,9\%\\
\bottomrule
\end{tabular}
\caption{Entrada de facção em 2025 entre os municípios sem facção em 2024, segundo o
garimpo em 2021--23.}\label{tab:garimpo}
\end{table}

%=====================================================================
\section{Etapa 4 --- MLE, MAP, Naive Bayes e EM}

\paragraph{MLE e MAP.} Para $\theta=P(\text{facção})$ por estado, $\hat\theta_{MLE}=k/n$ e,
com prior Beta(2,2), $\hat\theta_{MAP}=(k+1)/(n+2)$. Onde há poucos municípios, o MLE é
extremo: no Acre ($k=n=22$), $\hat\theta_{MLE}=1$, o que daria probabilidade zero a um
município sem facção; o MAP dá 0,958. Em Roraima, 0,933 contra 0,882. Nos estados grandes
(PA, MA, MT) os dois coincidem até a segunda casa: o prior só pesa quando há pouca
evidência.

\paragraph{Naive Bayes do zero.} Com suavização de Laplace e os seis atributos da Etapa 1,
acurácia de 77,2\% e balanceada de 0,737 em 5 dobras, igual ao \texttt{CategoricalNB}.

\paragraph{EM com variável oculta.} O Cartografias classifica cada município em ``sem
registro'', ``uma facção (domínio)'' ou ``disputa''. Supomos que esse regime $Z$ é oculto e
tentamos recuperá-lo só a partir de $x=\log(1+\text{homicídios 2021--23})$, com uma mistura
de 3 gaussianas ajustada pelo EM (\texttt{em.py}). A log-verossimilhança cresce
monotonicamente e estabiliza em 5 iterações (Figura~\ref{fig:em}). O EM nunca vê o rótulo,
mas os grupos acompanham o Cartografias (Tabela~\ref{tab:em}): no grupo menos violento,
1,7\% estão em disputa; no mais violento, 29\%.

\begin{figure}[t]
\centering
\fig{figuras/fig4_em.pdf}{0.98}
\caption{EM: log-verossimilhança por iteração e mistura estimada.}\label{fig:em}
\end{figure}

\begin{table}[t]
\centering\small
\begin{tabular}{@{}lcccccc@{}}
\toprule
Grupo & $\theta$ & hom.\ típicos & sem reg. & uma & disputa\\
\midrule
$Z=0$ & 0,31 & 3,6 & 0,887 & 0,096 & 0,017\\
$Z=1$ & 0,36 & 9,2 & 0,724 & 0,219 & 0,057\\
$Z=2$ & 0,33 & 23,1 & 0,342 & 0,370 & 0,288\\
\bottomrule
\end{tabular}
\caption{Grupos do EM e a composição segundo o Cartografias (linhas somam 1).}\label{tab:em}
\end{table}

Com a \emph{taxa} por 100 mil no lugar da contagem, a correspondência é mais fraca (o grupo
mais violento tem 15\% em disputa), e o EM leva 38 iterações. A facção está ligada ao
\emph{volume} de violência, que acompanha o tamanho do mercado, e não só à intensidade
relativa. Com a taxa, uma componente tenta colapsar sobre os municípios sem nenhum
homicídio; um piso de variância ($\sigma^2\ge0{,}05$) evita a singularidade.

\paragraph{NB puro $\times$ NB + EM.} Trocando os tercis de \texttt{TXH} pelo grupo $Z$
(ajustado só no treino de cada dobra), o Naive Bayes passa de 77,2\% para 77,7\%
(balanceada de 0,737 para 0,745). O ganho é pequeno: $Z$ traz o volume de violência, mas
boa parte dessa informação já estava em \texttt{PORTE}.

%=====================================================================
\section{Etapa 5 --- Rede bayesiana}

\paragraph{Estrutura.} O DAG (Figura~\ref{fig:dag}) tem 6 nós. \texttt{PORTE}$\to$\F:
cidades maiores concentram mercado consumidor, portos, rodovias e presídios \cite{r10}.
\texttt{FRONT}$\to$\F: as fronteiras e suas cidades são nós da rede da cocaína \cite{r11}.
\F$\to$\{\texttt{TXH}, \texttt{ESC}, \texttt{TXT}\}: disputas e execuções elevam
homicídios, tentativas e mortes sem causa esclarecida. A rede é usada no sentido
diagnóstico: observamos porte, estado e sintomas e inferimos $P(\F\mid\text{evidências})$.

\begin{figure}[t]
\centering
\fig{figuras/fig5_dag.pdf}{0.8}
\caption{DAG: \F{} é colisor de \texttt{PORTE} e \texttt{FRONT} e causa comum dos três
sintomas.}\label{fig:dag}
\end{figure}

\paragraph{CPTs e economia.} As CPTs são estimadas por máxima verossimilhança. A tabela
central, $P(\F{=}1\mid\texttt{PORTE},\texttt{FRONT})$, vai de 0,013 (pequeno, sem fronteira)
a 0,861 (grande, com fronteira). A conjunta das 6 variáveis ($3\cdot2\cdot2\cdot3\cdot2\cdot3=216$
células) exige 215 parâmetros; a rede, 19: $2+1+6+4+2+4$. Economia de 196 parâmetros (91\%).

\paragraph{Inferência por enumeração.} Implementamos
$P(X\mid e)=\alpha\sum_y\prod_i P(x_i\mid pais(X_i))$ do zero e comparamos com a
\texttt{VariableElimination} do \texttt{pgmpy} em 7 consultas (15 probabilidades): a maior
diferença é $1{,}1\times10^{-16}$. Exemplos: $P(F\mid\text{grande})=0{,}79$;
$P(F\mid\text{pequeno})=0{,}13$; um município pequeno de fronteira com os três sintomas no
nível mais alto tem $P(F)=0{,}43$; um grande, sem fronteira e sem sintomas, 0,29.

\paragraph{d-separação.} (1) \emph{Garfo}: \texttt{TXH}$\leftarrow$\F$\rightarrow$\texttt{ESC};
dado \F, os sintomas são d-separados. (2) \emph{Colisor}:
\texttt{PORTE}$\rightarrow$\F$\leftarrow$\texttt{FRONT}. Sem observar \F,
$P(\texttt{FRONT}\mid\text{pequeno})=P(\texttt{FRONT}\mid\text{grande})=0{,}559$:
independentes. Observando $\F{=}1$, o caminho se abre:
$P(\texttt{FRONT}\mid\text{pequeno},F)=0{,}955$ contra 0,608 para grande. É o
\emph{explaining away}: uma facção num município pequeno é mais bem ``explicada'' pela rota
do tráfico do que pelo tamanho. Nos dados, porte e fronteira têm correlação 0,21; a
independência marginal suposta pelo DAG é uma aproximação.

\paragraph{A rede como classificador.} Em 5 dobras, AUC $=0{,}81$ (acurácia 76,1\%,
balanceada 0,713). As probabilidades passadas à Etapa 6 são \emph{fora da dobra}: cada
município recebe a $P(F)$ de uma rede treinada sem ele, para que a validação seguinte não
seja otimista.

\paragraph{Deixando um estado de fora.} A AUC geral cai pouco (0,80), mas a tabela por
estado mostra o que a média esconde (Tabela~\ref{tab:uf}). A rede ordena bem os municípios
em TO, AP, MA e PA, mas erra o \emph{nível} onde a facção é quase universal: no Acre e em
Roraima a $P$ média prevista é cerca de 0,47, porque \texttt{FRONT} só distingue MA e TO dos
demais. Serve para ordenar municípios num estado conhecido, não para estimar o nível num
estado novo.

\begin{table}[t]
\centering\small
\begin{tabular}{@{}lcccc@{}}
\toprule
UF & $n$ & facção & $P$ média & AUC\\
\midrule
AC & 22 & 1,00 & 0,47 & --\\
RR & 15 & 0,93 & 0,47 & 0,64\\
PA & 144 & 0,51 & 0,53 & 0,78\\
RO & 52 & 0,50 & 0,43 & 0,58\\
AM & 62 & 0,32 & 0,32 & 0,65\\
AP & 16 & 0,31 & 0,51 & 0,85\\
MT & 141 & 0,30 & 0,51 & 0,75\\
MA & 217 & 0,22 & 0,21 & 0,79\\
TO & 139 & 0,06 & 0,11 & 0,94\\
\bottomrule
\end{tabular}
\caption{Validação deixando um estado de fora.}\label{tab:uf}
\end{table}

\paragraph{Sensibilidade ao rótulo.} Sem os 45 municípios que ``saíram'' da lista em 2025,
a AUC sobe para 0,83; a rede já os via com probabilidade intermediária (0,40, contra 0,58
dos que têm facção nas duas edições e 0,23 dos sem facção), o que é coerente com erro de
registro ou mudança de método. Com o rótulo ``facção em 2024 ou 2025'', a AUC é 0,82. As
conclusões não dependem do ruído do rótulo.

%=====================================================================
\section{Etapa 6 --- Lógica fuzzy e integração}\label{sec:fuzzy}

\paragraph{Variáveis linguísticas.} O sistema Mamdani (\texttt{scikit-fuzzy}) tem duas
entradas: \emph{presença} $=P(F)$ vinda da rede (baixa, média, alta, centradas em torno da
taxa-base de 32\%) e \emph{tendência} = variação relativa dos homicídios de 2020--21 para
2022--23 (queda, estável, alta). A saída é a \emph{prioridade} de monitoramento, de 0 a 100
(Figura~\ref{fig:pert}). Todas as funções de pertinência são triangulares, que em sistemas
Mamdani para nível de criminalidade classificaram melhor do que as trapezoidais \cite{r19}.

\begin{figure}[t]
\centering
\fig{figuras/fig6_pertinencia.pdf}{0.98}
\caption{Funções de pertinência.}\label{fig:pert}
\end{figure}

\paragraph{Regras.}
\begin{enumerate}
\item SE presença alta E (tendência alta OU estável) ENTÃO prioridade alta.
\item SE presença alta E tendência em queda ENTÃO prioridade média (possível domínio
consolidado).
\item SE presença média E tendência alta ENTÃO prioridade alta (possível expansão).
\item SE presença média E (estável OU queda) ENTÃO prioridade média.
\item SE presença baixa ENTÃO prioridade baixa.
\end{enumerate}
A regra 2 traduz um padrão descrito na literatura: onde uma facção já domina sozinha, a
violência pode cair (``paz do domínio''). A defuzzificação é pelo centroide; classificamos a
saída em baixa ($<40$), média (40--60) e alta ($\ge60$).

\paragraph{Conferência com 2024 e 2025.} A Tabela~\ref{tab:fuzzy} confere a decisão com o
rótulo de 2024 e com as entradas de 2025. A proporção com facção cresce de 13\% para 61\% da
prioridade baixa para a alta, e a de disputa, de 1,7\% para 25,5\%. Entre os 549 sem facção
em 2024, a taxa de entrada em 2025 vai de 13\% na baixa para 44\% na alta. É uma validação
fora do tempo: nenhum dado de 2025 foi usado.

\begin{table}[t]
\centering\small
\begin{tabular}{@{}lcccc@{}}
\toprule
Prioridade & Munic. & Facção 2024 & Disputa & Entrada 2025$^*$\\
\midrule
Baixa & 404 & 12,9\% & 1,7\% & 13,1\%\\
Média & 165 & 37,6\% & 10,3\% & 40,8\%\\
Alta & 239 & 60,7\% & 25,5\% & 43,6\%\\
\bottomrule
\end{tabular}
\caption{$^*$Entre os municípios sem facção em 2024. AUC para prever a entrada: $P(F)$
bayesiana 0,764; prioridade fuzzy 0,751.}\label{tab:fuzzy}
\end{table}

\paragraph{Fuzzy $\times$ limiar rígido.} O limiar $P\ge0{,}5$ divide os municípios em 234
``alta'' e 574 ``baixa''. Ele também separa bem (entrada em 2025: 54\% contra 18\%), mas
joga numa só classe os 465 que ficaram abaixo do limiar. O fuzzy cria uma faixa
intermediária de 165 municípios que já concentra 41\% de entradas, e a saída é contínua
perto da fronteira de decisão (Figura~\ref{fig:sup}): um município com $P=0{,}49$ e outro
com 0,51 recebem prioridades próximas, não opostas. Isso importa aqui porque o rótulo é
ruidoso (``sem registro'' não é ausência certa) e a fronteira entre as classes é ambígua:
são exatamente os casos em que o fuzzy supera o limiar rígido. A tendência também desloca a
prioridade, o que um limiar sobre $P$ não faz.

\begin{figure}[t]
\centering
\fig{figuras/fig7_fuzzy_limiar.pdf}{0.98}
\caption{Prioridade fuzzy em função de $P(F)$ para três tendências, contra o limiar
rígido.}\label{fig:sup}
\end{figure}

\paragraph{A violência de 2024 antecipa a chegada?} Com o VDE de 2024, testamos duas
hipóteses com uma regressão logística em 5 dobras, só entre os sem facção em 2024. (i)
Trocar a janela 2021--23 por 2022--24 não melhora (AUC 0,754 nas duas). (ii) Municípios
cujos homicídios saltaram em 2024 entraram mais (31\% contra 21\%), mas o salto de 2023 tem
o mesmo efeito (32\% contra 21\%). O salto é sinal de município violento, não da chegada da
facção. Estudos de caso, como o de Iranduba (AM), em que os homicídios subiram 318\% em 2021
após a criação de uma facção \cite{r15}, não aparecem como padrão na escala dos 808
municípios.

\paragraph{Interpretabilidade.} Cada regra é uma frase que um analista de inteligência
entende e pode contestar, e cada entrada tem origem rastreável: a presença vem da rede
(cujas CPTs são tabelas de frequência) e a tendência vem direto do VDE.

%=====================================================================
\section{Experimentos com outras variáveis}\label{sec:exp}

O enunciado limita a árvore a 4--6 atributos e a rede a 4--6 nós, e o pipeline já usa 6.
Testamos então 15 variáveis candidatas \emph{trocando} um nó da rede (em geral \texttt{ESC}),
com as mesmas validações (\texttt{experimentos\_variaveis.ipynb}); a
Tabela~\ref{tab:var} resume.

\begin{table}[t]
\centering\footnotesize
\begin{tabular}{@{}lccc@{}}
\toprule
Variável na rede & 2024 & Estado fora & Entrada 2025\\
\midrule
Atual & 0,810 & 0,800 & \textbf{0,764}\\
Faixa de fronteira no lugar de \texttt{FRONT} & \textbf{0,837} & \textbf{0,811} & 0,708\\
Vizinho com facção (malha IBGE) & 0,819 & 0,800 & 0,742\\
Urbanização (Censo 2022) & 0,819 & 0,788 & 0,731\\
Hierarquia urbana (REGIC) & 0,807 & 0,787 & 0,751\\
Latrocínio, feminicídio, mandados & 0,81--0,82 & 0,80--0,81 & 0,73--0,76\\
Suicídio, trânsito (placebos) & 0,813--0,815 & 0,80 & 0,752\\
Garimpo IBAMA / INPE & 0,813 / 0,809 & 0,79--0,80 & 0,747\\
Desmatamento (DETER) & 0,808 & 0,796 & 0,733\\
\bottomrule
\end{tabular}
\caption{AUC da rede com cada variável candidata (probabilidades fora da
dobra).}\label{tab:var}
\end{table}

\paragraph{Fronteira por município.} Trocar \texttt{FRONT} por estado pela faixa de
fronteira do IBGE melhora o diagnóstico (AUC 0,84) e corrige boa parte do erro de nível: com
o estado de fora, a $P$ média prevista sobe para 0,68 no Acre e 0,70 em Roraima. Mas a
previsão da entrada em 2025 cai para 0,71: nos estados de fronteira, a expansão foi tão
frequente no interior (38\%) quanto na faixa (36\%), e a variável municipal rebaixa
justamente o interior do Pará, do Mato Grosso e do Amazonas. Como o objetivo final é
antecipar a expansão, mantivemos \texttt{FRONT} por estado.

\paragraph{As demais.} Nenhuma muda a AUC de 2024 em mais de 0,01, e todas pioram a entrada
em 2025. Os placebos (suicídio, trânsito) rendem tanto quanto as variáveis de violência: os
sintomas da rede carregam sobretudo o tamanho e a urbanização do município. Vizinhos,
urbanização e hierarquia urbana repetem o que porte e estado já dizem. Numa logística só
para a entrada em 2025, nenhuma variável eleva a AUC (0,779) além do desvio entre dobras
(0,05). O garimpo dobra a taxa de entrada, mas atinge só cerca de 40 dos 549 municípios,
pouco para mudar a AUC do conjunto.

%=====================================================================
\section{Discussão}

A Tabela~\ref{tab:ref} compara os resultados com as faixas de referência do enunciado.
Ficamos dentro ou acima nas Etapas 1 e 2 (regressão). Ficamos abaixo na logística e no
NB + EM. A razão é o próprio problema: a taxa-base é 32\%, os atributos públicos (violência
e população) só capturam parte do fenômeno e o rótulo só tem positivos confiáveis
(Seção~\ref{sec:concl}). A AUC de 0,81 mostra que os modelos ordenam bem os municípios; o
que limita a acurácia é a sobreposição real entre municípios violentos sem facção e facções
que dominam sem violência. O EM converge em 5 iterações (contagem) ou 38 (taxa), fora da
faixa de 15--30, porque a tolerância é relativa ($10^{-4}$); a curva cresce
monotonicamente, como esperado.

\begin{table}[t]
\centering\small
\begin{tabular}{@{}lcc@{}}
\toprule
Métrica & Esperado & Obtido\\
\midrule
Árvore: treino sem poda & 85--92\% & 87,4\%\\
Árvore: teste 5-fold & 72--80\% & 78,6\%\\
Profundidade ótima & 3--5 & 3--5\\
$R^2$ simples & 0,55--0,65 & 0,672\\
$R^2$ multivariada & 0,70--0,80 & 0,763\\
$R^2$ Ridge & 0,72--0,82 & 0,762\\
Logística & 78--85\% & 75,9\%\\
Naive Bayes puro & 72--78\% & 77,2\%\\
Naive Bayes + EM & 80--86\% & 77,7\%\\
Enumeração = \texttt{pgmpy} & idêntico & dif.\ $10^{-16}$\\
\bottomrule
\end{tabular}
\caption{Resultados obtidos $\times$ faixas de referência do enunciado.}\label{tab:ref}
\end{table}

O fio condutor do pipeline é a passagem da certeza à incerteza graduada. A árvore produz
regras legíveis, mas a monotonicidade mostra que elas quebram com um fato novo. A
probabilidade substitui ``é/não é'' por graus de crença; o EM mostra que a variável oculta
deixa rastro nos dados; a rede organiza essas crenças com poucos parâmetros e permite
raciocinar nos dois sentidos; e o fuzzy devolve uma decisão legível, agora sem fronteiras
rígidas.

%=====================================================================
\section{Comunicação}

Houve comunicação com a equipe 5 sobre o trabalho:

\medskip
\noindent\textbf{Equipe 5 --- Segurança Pública:} Nosso projeto utiliza indicadores oficiais
de segurança pública para aplicar técnicas de Inteligência Artificial, incluindo árvores de
decisão, regressão, inferência probabilística, estimação estatística, redes Bayesianas e
lógica fuzzy. Buscamos identificar situações de elevado volume de ocorrências que possam
demandar maior atenção da gestão pública.

\smallskip
\noindent\textbf{Equipe 1 --- Facções Criminosas no Amazonas:} Nosso estudo investiga a
presença e a expansão de facções criminosas nos municípios do Amazonas. Utilizamos rótulos
provenientes das Cartografias, 3ª e 4ª edições, e atributos extraídos do Banco VDE do
SINESP/MJSP, abrangendo o período de 2020 a 2024. Também integramos informações
populacionais do IBGE e indicadores externos do IBAMA, INPE e IBGE.

\smallskip
\noindent\textbf{Equipe 5 --- Segurança Pública:} Uma diferença importante entre os projetos
está na unidade de análise. Trabalhamos com indicadores agregados por unidade federativa e
período, enquanto vocês utilizam informações municipais. A granularidade municipal pode
permitir análises territoriais mais específicas, embora exija cuidados adicionais com a
qualidade e a integração dos dados.

\smallskip
\noindent\textbf{Equipe 1 --- Facções Criminosas no Amazonas:} Exatamente. Nossa metodologia
utiliza códigos municipais para integrar diferentes fontes. Além dos registros de violência,
consideramos variáveis territoriais e ambientais que podem contribuir para investigar
associações com a presença de facções. Essas associações, entretanto, não devem ser
interpretadas automaticamente como relações causais.

\smallskip
\noindent\textbf{Equipe 5 --- Segurança Pública:} Essa abordagem oferece uma perspectiva
complementar à nossa. As técnicas probabilísticas e de classificação utilizadas em nosso
trabalho poderiam, em estudos futuros, ser aplicadas a indicadores municipais para
investigar padrões territoriais, desde que fossem realizados novos treinamentos e
validações.

\smallskip
\noindent\textbf{Equipe 1 --- Facções Criminosas no Amazonas:} Da mesma forma, a utilização
de redes Bayesianas e lógica fuzzy poderia contribuir para representar incertezas e produzir
classificações interpretáveis. Entretanto, seria necessário avaliar a qualidade dos rótulos
e os riscos de generalização indevida dos resultados.

\smallskip
\noindent\textbf{Equipe 5 --- Segurança Pública:} Concluímos que os projetos são
complementares: enquanto nossa pesquisa analisa indicadores agregados de segurança pública
para apoiar o monitoramento, o estudo de vocês investiga padrões municipais relacionados à
presença e à expansão de facções criminosas. A integração dessas perspectivas poderia
fundamentar pesquisas futuras sobre análise territorial e apoio à decisão pública.

\smallskip
\noindent\textbf{Síntese:} A comparação evidencia a importância da granularidade
territorial, da integração de bases oficiais, da interpretação probabilística e da validação
dos modelos. Também demonstra que indicadores de violência e presença de organizações
criminosas representam fenômenos distintos, que não devem ser tratados como equivalentes.

%=====================================================================
\section{Conclusão e limitações}\label{sec:concl}

O porte do município e o estado explicam a maior parte da presença registrada de facção na
Amazônia Legal; a violência refina essa estimativa. A probabilidade calculada com dados até
2023 antecipa parte da expansão observada em 2025 (AUC 0,76), mas a violência indica onde a
facção tende a estar, não quando ela chega. O garimpo, medido pelo IBAMA e pelo INPE, é a
pista mais clara de para onde ela vai: dobra a taxa de entrada em 2025. Das 15 variáveis de
outras bases que testamos, nenhuma melhorou os modelos de forma consistente; o limite parece
estar no rótulo, não na falta de atributos.

\paragraph{Limitações.} (1) ``Sem registro'' no Cartografias significa facção não
identificada pelos pesquisadores, não ausência certa; por isso falamos em facção registrada.
Formalmente, é aprendizado com positivos e não rotulados \cite{r18}: se os registros fossem
uma amostra aleatória dos municípios com facção, a probabilidade estimada diferiria da
verdadeira por um fator constante. A ordenação (e a AUC) se preservaria, mas o nível e a
acurácia ficariam subestimados. (2) Parte das mudanças entre a 3ª e a 4ª edição vem da
melhoria do método; 45 municípios ``saíram'' da lista. (3) O relatório tem pequenas
divergências entre as tabelas por estado e o quadro-resumo (AP, PA, TO); definimos
``disputa'' pelo número de facções listadas. (4) O Maranhão entra inteiro (217 municípios),
mas só 181 pertencem à Amazônia Legal. (5) \texttt{FRONT} é definido por estado; a versão
por município corrige o erro de nível em AC e RR, mas piora a previsão de 2025
(Seção~\ref{sec:exp}). (6) Não há estimativa do IBGE para 2023; as taxas usam o Censo 2022.

\paragraph{Ética.} A saída é uma prioridade estatística para orientar monitoramento, não uma
acusação. Modelos treinados com registros de fiscalização tendem a prever onde há
fiscalização, e não só onde há crime \cite{r20}; isso vale para o Cartografias e para os
autos do IBAMA. Municípios com prioridade alta e sem registro (por exemplo, Oriximiná, PA,
com $P=0{,}92$) são candidatos a verificação, não municípios ``dominados'': como nos
``falsos positivos'' de \cite{r17}, podem ser municípios com facção ainda não registrada.
Rotular publicamente uma localidade pode gerar estigma e afetar investimento e turismo; por
isso a lista de candidatos fica restrita ao notebook.

%=====================================================================
\begin{thebibliography}{99}\small
\setlength{\itemsep}{0pt}
\bibitem{r1} S. Russell e P. Norvig. \emph{Artificial Intelligence: A Modern Approach}, 4.\ ed.
Pearson, 2020 (caps.\ 18--20).
\bibitem{r2} L. A. Zadeh. Fuzzy sets. \emph{Information and Control}, 8(3):338--353, 1965.
\bibitem{r3} E. Hüllermeier. Fuzzy machine learning and data mining. \emph{WIREs DMKD},
1(4):269--283, 2011.
\bibitem{r4} Instituto Mãe Crioula e Fórum Brasileiro de Segurança Pública.
\emph{Cartografias da Violência na Amazônia}, 3.\ ed., 2024.
\bibitem{r5} Instituto Mãe Crioula e Fórum Brasileiro de Segurança Pública.
\emph{Cartografias da Violência na Amazônia}, 4.\ ed., 2025.
\bibitem{r6} Ministério da Justiça e Segurança Pública. Banco VDE (SINESP), 2020--2024.
\url{https://dadosabertos.mj.gov.br}. Acesso em out.\ 2026.
\bibitem{r7} IBGE. Censo Demográfico 2022, Estimativas da População, Municípios da Faixa de
Fronteira 2024, Malha Municipal e REGIC 2018. \url{https://www.ibge.gov.br}. Acesso em nov.\
2026.
\bibitem{r8} IBAMA. Dados abertos: autos de infração e termos de destruição.
\url{https://dadosabertos.ibama.gov.br}. Acesso em nov.\ 2026.
\bibitem{r9} INPE. DETER Amazônia, TerraBrasilis. \url{https://terrabrasilis.dpi.inpe.br}.
Acesso em nov.\ 2026.
\bibitem{r10} Í. B. L. Siqueira e L. F. S. Paiva. ``No Norte tem Comando'': as maneiras de
fazer o crime, a guerra e o domínio das prisões do Amazonas. \emph{Revista Brasileira de
Sociologia}, 7(17):125--154, 2019. DOI 10.20336/rbs.486.
\bibitem{r11} A. C. O. Couto. Fronteiras e estrutura espacial do narcotráfico na Amazônia.
\emph{Boletim Gaúcho de Geografia}, 47(1), 2020.
\bibitem{r12} A. C. O. Couto. Relações transfronteiriças do narcotráfico na Amazônia: dos
crimes conexos aos desafios da segurança regional. \emph{Boletim de Análise
Político-Institucional (Ipea)}, 36, 2024. DOI 10.38116/bapi36art5.
\bibitem{r13} R. P. Chagas. O ``narcogarimpo'' na Terra Indígena Yanomami. \emph{Boletim de
Análise Político-Institucional (Ipea)}, 36, 2024. DOI 10.38116/bapi36art7.
\bibitem{r14} R. Pereira-Chagas. Narcogarimpo: las afinidades electivas entre los frentes de
minería ilegal y la expansión del narcotráfico en la Amazonía brasileña. \emph{URVIO},
38:32--48, 2024. DOI 10.17141/urvio.38.2024.6144.
\bibitem{r15} P. V. A. Sales, L. C. Sena, F. C. Bezerra Neto e C. M. A. Mello. O impacto da
criação da facção Revolucionários do Amazonas (RDA) nos índices de violência letal em
Iranduba-AM no período de 2017 a 2024. \emph{Revista Geopolítica Transfronteiriça},
9(4):1--15, 2025.
\bibitem{r16} J. Ribeiro, L. Meneses, D. Costa, W. Miranda e R. Alves. Prediction of Homicides
in Urban Centers: A Machine Learning Approach. \emph{IntelliSys}, 2021. arXiv:2008.06979.
\bibitem{r17} G. M. Campedelli, G. Daniele e M. Le Moglie. Mafia, Politics and Machine
Predictions. CEPR Discussion Paper 19322, 2024.
\bibitem{r18} C. Elkan e K. Noto. Learning Classifiers from Only Positive and Unlabeled Data.
\emph{KDD}, p.\ 213--220, 2008.
\bibitem{r19} N. M. Mohamad Noor e S. H. Ab Hamid. Comparison of Triangular and Trapezoidal
Membership Functions for Improving Decision Making in Crime Prevention. \emph{Frontiers in
Artificial Intelligence and Applications}, 261, 2014. DOI 10.3233/978-1-61499-399-5-69.
\bibitem{r20} K. Lum e W. Isaac. To Predict and Serve? \emph{Significance}, 13(5):14--19,
2016.
\end{thebibliography}

%=====================================================================
\appendix
\section{Uso de assistente de IA}

Usamos um assistente de IA (Claude, via Claude Code) como apoio à implementação do código e
à análise estatística, e para revisar parte do texto. As decisões do trabalho (a pergunta de
pesquisa, a escolha dos atributos, o protocolo de validação e a interpretação dos
resultados) foram da equipe. Todo o código foi lido, executado e conferido por nós, e os
números do relatório vêm da execução dos notebooks com \texttt{seed = 42}.

\paragraph{Preparação dos dados.} O assistente ajudou a escrever os scripts que montam a base
municipal do VDE (2020--2024) e a juntam à população do IBGE pelo código do município, a
extrair do PDF do Cartografias 2025 os quadros de facções por estado
(\texttt{extrair\_cartografias\_2025.py}) e a agregar por município as bases externas usadas
na Seção~\ref{sec:exp} (\texttt{montar\_bases\_externas.py}).

\paragraph{Implementação dos algoritmos.} Com o assistente, implementamos do zero o ID3, a
regressão com gradiente descendente, a regressão logística, o Naive Bayes, o EM para mistura
de gaussianas e a inferência por enumeração na rede bayesiana. Cada implementação foi
comparada com a biblioteca correspondente (scikit-learn e \texttt{pgmpy}); a coincidência
dos resultados, por exemplo a diferença máxima de $10^{-16}$ na enumeração, serviu de
verificação do código.

\paragraph{Análise estatística.} O assistente auxiliou a tornar a validação mais rigorosa:
tercis calculados apenas no treino de cada dobra, probabilidades da rede obtidas fora da
dobra, validação deixando um estado de fora e validação fora do tempo com a 4ª edição do
Cartografias. Também ajudou a organizar os experimentos com as 15 variáveis candidatas
(\texttt{experimentos\_variaveis.ipynb}), os testes de hipótese e as tabelas comparativas do
relatório.

\paragraph{Revisão bibliográfica e entrega.} O assistente sugeriu artigos relacionados ao
tema; títulos, autores e trechos citados foram conferidos por nós nos PDFs. Ele também
ajudou a preparar o relatório em \LaTeX, o README e os notebooks para a entrega.

\section{Reprodução}

\texttt{pip install -r requirements.txt}; \texttt{python utils.py}; executar
\texttt{etapa1.ipynb} a \texttt{etapa6.ipynb} em ordem (a Etapa 5 grava
\texttt{dados/probabilidades\_bayesianas.csv}, lido pela Etapa 6). Repositório:
\url{https://github.com/RomarioAndrade/seguranca-publica-ml-fuzzy}.

\end{document}

